%-----------------------------------------------------------------------------%
\section{Conclusion and Future Work}
\label{bab:5}
%-----------------------------------------------------------------------------%

This research designed and evaluated \textit{ResourceDefragmentation}, a request-based Kubernetes Descheduler \textit{Balance} policy that repairs post-placement CPU-memory fragmentation without replacing the kube-scheduler.

\textbf{Identifying fragmented Nodes.} \textit{ResourceDefragmentation} identifies fragmented Nodes directly from their request-space resource shape, without relying on runtime usage or a single utilization threshold. It selects drain candidates from average utilization and CPU-memory bin quality, orders them by the Resource Imbalance Index (RII) and Free-Space Index (FSI), and then evicts the Pod whose predicted landing Node, found through \textit{predictSchedulerTarget}, has the best projected bin score. A target ceiling and pack-upward guard prevent overfilling a Node or relocating a workload to an emptier one. Together the drain ordering, predicted-target criterion, and guards drive effective packing while keeping eviction churn low.

\textbf{Effectiveness against the baseline.} ResourceDefragmentation reclaimed four Nodes in S1, two in S2, and one in S4, all cases where \textit{HighNodeUtilization} performed no eviction. It cut dimensional stranding by \(73\%\) in S2 (\(3.495 \to 0.961\)) and \(79\%\) in S4 (\(3.071 \to 0.633\)) in the layouts where stranding is reducible, and on the already-irreducible layouts (S3, S5) it correctly declined unproductive evictions rather than churning. Four of five scenarios converged within three passes and no run left a Pod Pending.

\textbf{Component attribution.} An ablation that disabled one mechanism at a time showed that each earns its place on a different axis: RII/FSI ordering preserves stranding quality (removing it regressed S4 from \(0.633\) to \(1.581\)), the pack-upward guard bounds eviction churn (removing it roughly doubled evictions for no stranding gain), and the projected-bin-score selector lowers the eviction count without changing the converged stranding. The stranding wins therefore come from the consolidation-and-defragmentation pipeline as a whole, while the selector keeps disruption low.

Overall, the results support a simple conclusion: selecting eviction candidates by the balance of their predicted landing Nodes, combined with request feasibility, bounded utilization, drain ordering, and pack-upward behavior, is enough to drive useful Node reclamation and dimensional defragmentation where it is achievable while avoiding wasteful eviction where it is not.

\subsection{Future Work}
\label{sec:future_works}

\begin{itemize}
    \item \textbf{Broader Resource Dimensions:} Extend the framework to additional resource dimensions, including storage and GPU, so that target prediction and bin-quality scoring can handle workloads beyond CPU and memory.
    \item \textbf{Scheduler Coordination:} Establish and evaluate scheduler coordination mechanisms that expose the predicted target, or a compatible scoring signal, to the kube-scheduler, ensuring that predicted-target placements align with actual scheduling decisions and eliminating the misalignment that currently bounds maximum defragmentation precision.
    \item \textbf{Pending-Pod Awareness:} Introduce Pending-pod awareness into the eviction pipeline so that drain candidates are selected with knowledge of already-unschedulable workloads, ensuring evictions free the right resource type on the right Node to unblock Pending Pods rather than delaying them further.
    \item \textbf{Threshold Sensitivity Analysis:} Conduct a sensitivity analysis of the drain threshold, target ceiling, and related parameters across various workloads and cluster types to characterize their effect on node reclamation, stranding reduction, and convergence stability.
    \item \textbf{Broader Generalizability Evaluation:} Evaluate the plugin across broader cluster shapes, node heterogeneity, and workload mixes to assess generalizability beyond the five scenarios studied.
    \item \textbf{Broader Baselines and Multi-Seed Runs:} Extend the comparison beyond \textit{HighNodeUtilization} once faithful reference implementations of prior pseudocode-only proposals become available, and repeat each policy-scenario cell across multiple seeds to quantify run-to-run variation.
\end{itemize}
